\section{Liquidity and subsequent opportunities}
\subsection{Cash accompanies engagement and repricing}
Low cash often accompanies bargaining engagement rather than immediate acceptance. Among all proposer-only completion requests, recipients below \$100 reject 17/69, counter 44/69 and accept 8/69. At \$300 or more they reject 120/175, counter 48/175 and accept 7/175. Stage and configuration contribute to the contrast. Forty-three low-cash requests are later offers, and 21 face Opus. Within opening requests to Sol, Gemini and Grok, the corresponding rejections are 13/20 and 105/125. Non-completing requests also receive fewer refusals below \$100, 4/9 versus 92/138. The association is not specific to conceding a rival's completion.

The pattern also differs across systems. Across opening and later completion requests, Sol rejects 4/26 below \$100 versus 25/45 at \$300 or more; Grok rejects 3/11 versus 46/52. Gemini remains resistant at 10/11 versus 46/52. Opus rejects none of 21 low-cash requests and only three of 26 high-cash ones. Cash on hand is not total financing capacity, and these counts do not establish a common distress policy.

Fourteen settled one-sided completing transfers have earlier resistance by the same holder to the same prospective completer. Pairing uses the first earlier refusal or counter on a preceding turn. Cash falls in thirteen pairs, from a median \$571.50 to \$82. Only five preserve the reciprocal deed bundle; eight improve net cash terms for the holder. Selection on settlement, changed offers and changed states prevent a cash-only explanation. The histories nevertheless show how a previously withheld deed enters a later financing decision.

Asking prices supply a similar qualified result. Completing-buyer median asks exceed other-buyer medians in six Sol/Opus/Grok-by-mortgage cells. Gemini's corresponding completing cell has one observation and does not enter that six-cell claim. Among six same-game/seller/deed overlaps retaining mortgage status, four are higher, one lower and one equal. The contrasts describe asks rather than an identified completion premium. Appendix incidents G2 and G7 connect pricing to the buyer's retained building cash and the seller's next construction.

\Needspace{20\baselineskip}
\subsection{Financing a blocker against a third party}
\TranscriptTitle{insurance}{An inheritance pitch selectively presents collateral}
G13 Grok has \$35, owns yellow and stands on Atlantic before Sol's Short Line. It offers mortgaged Pennsylvania to Opus for \$200:
\TranscriptQuote{g13-opening-insurance}
Opus counters at \$130. Grok holds at \$200 while its applied note recognizes the available \$130 Ventnor mortgage. Its provider explanation says:
\TranscriptQuote{g13-omission-explanation}
Opus accepts \$200 and pays \$16 transfer interest. Its returned explanation repeats the estate concern after the appeal, rather than establishing an independently recorded prior motive. The seller's stated minimum omits a resource: \$130 proceeds plus \$35 cash and the available mortgage cover the threatened \$200 rent. The stored conditional check remains a historical witness; no alternative route is executed here.

A direct dice move from Atlantic to Short Line requires nine, with probability 4/36. The public 5/36 is wrong for that event; card movement and extra rolls require separate accounting. Grok misses Short Line next roll but later goes bankrupt there at turn 150, transferring yellow to Sol. Opus completes green at auction seven turns after the purchase. The episode supports selective presentation of financing and receipt of an estate-related pitch. It does not identify whether that omission caused the extra \$70 payment or quantify the survival purchased.
\TranscriptSource{G13 d250--253 and d311. Public messages occur in the actual recipient requests. Opus uses action-only recovery. The omission text is Grok's provider-returned explanation, not a message delivered to Opus.}

\subsection{Groups become productive, then change hands}
After a trade completes a group, Sol and Gemini build in the first management window on 16/17 and 17/18 groups; Opus and Grok do so on 7/15 and 3/13. The twenty other windows include nine without a build option for that group, four defaults, four content/call conflicts and three valid ends. A missing immediate build does not itself establish that no financing sequence could enable construction. Trade creates an investment option, not an automatic income stream.

Trade-completed groups receive \$57,378 before ownership loss or game end, with \$53,810 going to the 43 built immediately. Median remaining ownership durations are 62 and 63.5 turns, but property types, landing hazards, mortgage status and later construction remain unbalanced. Eleven other groups build later. These are realized receipts rather than incremental returns caused by the first-window choice.

Nonterminal estates transfer 227 deeds and complete 52 groups, against 207 deeds and 63 completions through trade. Thirty-eight of forty nonterminal eliminations follow rent on the creditor's developed color property. The first estate recipient wins eighteen games but already leads in buildings in twelve and cash in eight; those overlapping prior-strength measures cannot be added. A further 69 completions occur at final elimination with no subsequent earning opportunity. The possible creditor is therefore a substantive object of anticipation, while inheritance's independent contribution to winning remains unmeasured.
